% Encodings: decoding first loses, computing on compressed data wins, and the decode budget.

\begin{table}[t]
\centering\small
\caption{Compression ratio (uncompressed bytes / encoded bytes) per data set.}
\label{tab:encratio}
\gen{encodings_ratio}
\end{table}

Each data set has one encoding it is meant for (\cref{tab:encratio}): uniform values in
$[0,100)$ for bitpack, for and dict (7 bits, 9.1$\times$; dictb's 1-byte codes 8.0$\times$);
sorted values for delta (15.6$\times$) and for (5.3$\times$); runs of mean length 64 for rle
(32$\times$); Zipf values from 10{,}000 distinct ones for dict (4.5$\times$, dictb 4.0$\times$);
and \emph{wide} 64-bit values, where nothing compresses. The table has $N{=}8$ attributes and
4M rows, from DRAM.

\subsection{Decoding first never pays off}
\label{sec:enc:decode}

\begin{table*}[t]
\centering\small
\caption{Speedup over plain DSM, best strategy, scalar / AVX2 kernels. Above 1 the encoding is
faster than the uncompressed column.}
\label{tab:encspeed}
\resizebox{\textwidth}{!}{\gen{encodings_speedup}}
\end{table*}

The early strategy decodes every vector and then runs the plain kernel. With scalar code it is
slower than plain DSM for every encoding, whatever the compression: SUM of all 8 attributes
runs at 0.80$\times$ for bitpack (9.1$\times$ smaller), 0.69$\times$ for rle (32$\times$
smaller), 0.30$\times$ for delta (15.6$\times$) and 0.30$\times$ for dict. The speed does not
follow the compression ratio, because the scan is no longer limited by bytes but by the
instructions that decode them. Value mode, one virtual call per value, is 5--11$\times$ slower
than plain for every encoding, plain included.

\subsection{Computing on the compressed form}
\label{sec:enc:compressed}

\begin{table}[t]
\centering\small
\caption{Decoding first vs working on the compressed form (ms, AVX2 kernels). SUM rows: all 8
attributes, on the encoding's own data set. Zone maps and dict: SUM(c0) with a predicate at
0.1--1\%.}
\label{tab:encdecode}
\begin{tabular}{llrrr}
\toprule
Encoding & Operation on the compressed form & decode first & compressed & plain \\
\midrule
rle & $\sum$ value $\times$ run length & 11.5 & \textbf{0.10} & 8.4 \\
delta & $m f + \sum_k d_k (m-k)$ per block & 24.5 & \textbf{6.4} & 8.0 \\
bitpack & $m \cdot \min + \sum$ offsets & 7.7 & \textbf{4.9} & 7.9 \\
for & zone maps skip blocks (sorted) & 4.0 & \textbf{0.09--0.36} & 1.3--1.4 \\
dict & compare 14-bit codes (unpacked first) & 2.6 & 1.63 & \textbf{1.28} \\
\bottomrule
\end{tabular}
\end{table}

Without a predicate, late SUM on an encoded column calls one aggregate that works on the
compressed data, so the work scales with the compressed size, not the number of values
(\cref{tab:encdecode}).
\begin{itemize}
\item \textbf{rle} adds value $\times$ run length once per run: 0.10\ms, 87$\times$ faster than
plain DSM. The first version took 1.9\ms, because a binary search for the first run of every
1024-row vector took 60\% of the time (each vector covers only about 16 runs). Remembering the
run the previous vector ended in, one call per column when there is no predicate, and a loop
over whole runs that gcc vectorizes brought it to 0.10\ms. The 7.6\,MB of runs then stay in L3
between repetitions, which is part of the 87$\times$.
\item \textbf{delta} sums each block from its deltas, $m f + \sum_k d_k(m-k)$, without the
prefix sum that rebuilds the values. The prefix sum is a dependency chain through every value
(3--4 cycles per value); the formula has none, and the AVX2 kernel computes it with two
running sums per lane: 6.0--6.4\ms, 1.25--1.31$\times$ faster than plain.
\item \textbf{Predicates on packed data.} bitpack and for compare packed offsets against a
translated bound and extract the other attributes only at selected positions. With 10-bit
values a line holds 51 of them, so late gathers 6.4 times fewer lines, and late SUM on bitpack
beats every uncompressed plan up to 5\% (2.4\ms{} at 1\% against 2.7 for late and 3.6--3.9 for
early on plain DSM). The same packing lowers the crossover within one encoding, 5--10\% for
SUM on bitpack and for against 17--21\% for the uncompressed column, because $f(s)$ is already
0.40 at 1\% (\cref{sec:mat:threshold}).
\item \textbf{Zone maps.} On sorted data, for skips every block whose minimum and maximum
exclude the predicate and accepts every block they include: 0.09--0.36\ms{} at 0.1--1\%,
against 1.3--1.4\ms{} for late on plain DSM and 4.0\ms{} for decoding first.
\end{itemize}

\subsection{Dictionary codes that fit a SIMD lane}
\label{sec:enc:dictb}

\begin{table*}[t]
\centering\small
\caption{\code{SUM WHERE c0 < x} on Zipf data in ms: plain DSM, dict (14-bit packed codes) and
dictb (2-byte codes), best of early and late, AVX2. The selectivities are those the Zipf data
allows.}
\label{tab:encdict}
\gen{enc_dict}
\end{table*}

dict is the exception in \cref{tab:encdecode}: working on its codes still loses (1.63 vs
1.28\ms{} at 1\%), because each code is stored in exactly $\lceil\log_2 |D|\rceil = 14$ bits and
has to be unpacked into a 32-bit lane before the cheap comparison. dictb rounds the code width
up to 1, 2 or 4 bytes, a lane width AVX2 compares directly. A predicate is then one
subtraction, one unsigned minimum and one comparison per 16 two-byte codes read straight from
memory, and \code{movemask} gives one bit per code. It costs 12\% of the compression
(4.5$\times$ to 4.0$\times$) and turns the dictionary predicate from a loss into a win
(\cref{tab:encdict}): 0.57\ms{} at 1\%, 2.2$\times$ faster than plain and 2.9$\times$ faster
than dict, and 5.4$\times$ faster than plain at 0.05\%. The gain shrinks as more is gathered
(0.9$\times$ at 41\%, where every selected row still needs a lookup), and scalar code gets
nothing from the byte layout (1.9\ms{} at every selectivity).

\subsection{When decoding must happen: a budget per value}
\label{sec:enc:budget}

\begin{table}[t]
\centering\small
\caption{SUM of 8 attributes on uniform data, cycles per value, by where the uncompressed table
would live. Packed: bitpack and for with the AVX2 fused unpack-and-add.}
\label{tab:budget}
\begin{tabular}{lrrr}
\toprule
Table in & Plain & Packed (AVX2) & Packed vs plain \\
\midrule
L2 / L3 & 0.26--0.30 & 0.61--0.68 & 0.40--0.49$\times$ \\
16\,MB (spilling L3) & 0.46 & 0.61--0.68 & 0.72--0.76$\times$ \\
DRAM & 0.90--0.98 & 0.61--0.68 & \textbf{1.42--1.60$\times$} \\
\bottomrule
\end{tabular}
\end{table}

Where the values must be decoded, decoding and loading overlap, so the time per value is
\[
  t = \max(\text{compute},\ \text{load}),\qquad
  \text{packed wins when}\ \max(\text{decode}+\text{add},\ \text{load}_\text{packed})
  < \max(\text{add},\ \text{load}_\text{plain}).
\]
From DRAM the packed load is small (9$\times$ less data), so packed wins when decoding plus
adding costs less than loading a plain value: about 1 cycle per 8-byte value at 28\,GB/s. In
cache the plain load is cheap and the budget shrinks to 0.26--0.30 cycles. The measurements
follow it (\cref{tab:budget}).
\begin{itemize}
\item \textbf{Scalar unpacking} of 7-bit values costs 1.2--1.3 cycles per value, just over the
DRAM budget: bitpack moves 9.1$\times$ fewer bytes and is still 0.8$\times$ as fast as plain.
\item \textbf{AVX2 unpacking with a fused sum} costs 0.61--0.68 cycles per value at every size.
In the format we use, 8 consecutive $w$-bit values occupy exactly $w$ bytes; for $w \le 25$,
one byte shuffle places each value's bytes in its own 32-bit lane and a per-lane variable
shift (\code{vpsrlvd}, AVX2 only) and a mask extract it, and the offsets are added without
being stored. From DRAM that is within the budget, and bitpack and for become
1.56--1.60$\times$ faster than plain on uniform data; in cache it is twice the plain cost.
\item \textbf{Dictionary decoding} adds a dependent lookup per value (2--4 cycles), over the
budget even with SIMD unpacking. A dictionary pays off for predicates, not for decoding.
\end{itemize}
So on this machine compression is a way to make a DRAM-bound scan CPU-bound, and it pays off
only from DRAM and only with SIMD decoding: the deciding size is the L3 capacity, measured on
the uncompressed table.
